\chapter{Introduction}
\label{chap:intro}

\section{Background}

Automated facial analysis systems, covering perceived gender classification, age estimation, and race or ethnicity prediction, are now embedded across commercial and public-sector applications ranging from photo tagging and demographic analytics to surveillance and access control. As these systems have proliferated, a substantial body of empirical research has documented that many of them exhibit statistically significant accuracy disparities across demographic subgroups, particularly by skin tone, gender presentation, and age. The foundational audit by Buolamwini and Gebru demonstrated that leading commercial gender-classification services performed markedly worse on darker-skinned female subjects than on lighter-skinned male subjects, with error rates as high as 34.7\% for the worst-performing intersectional subgroup against as low as 0.8\% for the best-performing one \cite{buolamwini2018gendershades}. That finding catalysed a wider research agenda into fairness auditing for computer vision, subsequently surveyed comprehensively by Dehdashtian et al., who catalogue the sources of bias that arise across the machine learning pipeline, from data collection and annotation through model training and deployment, and the corresponding families of fairness metrics and mitigation techniques proposed in response \cite{dehdashtian2024survey}.

Despite this growing awareness, independent, systematic, and reproducible auditing of facial analysis models remains comparatively rare in practice. Most organisations that deploy such systems either do not audit them for demographic bias at all, or rely on ad hoc, non-standardised evaluations that are difficult to reproduce, compare, or scrutinise externally. Dehdashtian et al.'s survey confirms that while a rich taxonomy of fairness metrics and mitigation techniques exists in the computer vision literature, comparatively little attention has been paid to standardising the auditing workflow itself into a repeatable pipeline \cite{dehdashtian2024survey}.

This gap is now closing a regulatory window rather than an academic one. The \gls{eu} \gls{ai} Act (\gls{eu} Regulation 2024/1689) enters into force for Article~10 (Data and Data Governance) on 2~August~2026, imposing binding data-governance, representativeness, and bias-detection obligations on providers of high-risk \gls{ai} systems, a category that explicitly includes biometric categorisation \cite{euaiact2024}. The \gls{nist} \gls{ai} Risk Management Framework (\gls{rmf}) offers a complementary, voluntary structure, organised around four iterative functions, Govern, Map, Measure, and Manage, for operationalising exactly this kind of risk assessment \cite{nist2023aiframework}. Organisations that have not embedded documentation and measurement infrastructure into their development lifecycle ahead of this deadline face substantially higher retrofit costs than those that plan for it from the outset.

\section{Problem Statement}

There is therefore a clear engineering gap between a mature and growing body of fairness metrics, benchmark datasets, and bias-mitigation research, and the availability of an accessible, reusable, standards-aligned software platform that operationalises this research into a repeatable auditing workflow for facial analysis systems specifically. Existing audits, including the Gender Shades study itself, are typically conducted externally, manually, and for a single point in time; they are not designed as reusable, repeatable software artefacts that other researchers or organisations can apply to arbitrary models. This leaves model owners, auditors, and researchers without a systematic way to evaluate a facial-analysis model's performance across demographic subgroups using standardised benchmark datasets, statistically grounded fairness metrics, and a transparent, shareable report, at precisely the moment regulatory frameworks such as the \gls{eu} \gls{ai} Act begin to require exactly this kind of evidence.

\section{Objectives}

\subsection{General Objective}
To design and propose an end-to-end diagnostic and evaluative software platform, BiasAperture, that systematically identifies demographic accuracy disparities in third-party facial analysis models and reports them in a standardised, regulator-legible format.

\subsection{Specific Objectives}
\begin{enumerate}
    \item To design a modular software architecture that accepts a trained facial analysis model, or a file of its precomputed predictions, together with a demographically labelled benchmark dataset, and computes subgroup and intersectional fairness metrics from it.
    \item To integrate established open-source fairness toolkits, AIF360 and Fairlearn, alongside standard statistical significance testing, so that every disparity the platform reports is methodologically defensible rather than an artefact of sampling noise.
    \item To generate a standardised, exportable fairness report, structurally informed by the Model Cards \cite{mitchell2019modelcards} and Datasheets for Datasets \cite{gebru2018datasheets} documentation conventions, that communicates subgroup performance disparities in a form accessible to both technical and non-technical stakeholders.
    \item To validate the platform's design against publicly available, demographically annotated benchmark datasets; the current case study uses FairFace \cite{karkkainen2021fairface}, while UTKFace was profiled and cut from the implementation scope.
    \item To map every computed fairness metric to a specific obligation under Article~10 and Annex~IV of the \gls{eu} \gls{ai} Act and to the corresponding function of the \gls{nist} \gls{ai} \gls{rmf}, so the platform's output is directly usable as compliance evidence.
\end{enumerate}

\section{Scope and Limitations}

\subsection{Scope}
BiasAperture is scoped as a strictly diagnostic and evaluative platform. In scope are: (i) demographic subgroup and intersectional performance evaluation of a supplied facial-analysis classifier; (ii) computation of four disparity metrics, \gls{dpd}, \gls{eod}, \gls{eop}, and \gls{dir}, backed by chi-squared significance testing and bootstrap confidence intervals; (iii) explainability to help identify which visual features are associated with a detected disparity --- the current explainability implementation uses demographic-dummy surrogate attribution, and richer spatial \gls{shap} and ITA analysis remain deferred; and (iv) automated generation of a Model-Cards- and Datasheets-style \gls{html} fairness report, with each metric traceable to its \gls{eu} \gls{ai} Act and \gls{nist} \gls{ai} \gls{rmf} basis. The current case study uses FairFace \cite{karkkainen2021fairface}; UTKFace was profiled and cut from the implementation scope, with a FairFace-trained \gls{cnn} baseline serving as the minimum case study.

\subsection{Limitations}
<<<<<<< Updated upstream
Consistent with these boundaries, BiasAperture does not attempt to mitigate bias, retrain or fine-tune models, or generate synthetic demographic data; its function is strictly diagnostic, not corrective. It reports where disparities exist and their statistical strength, but it does not prescribe or apply any remediation. Its findings are only as reliable as the label quality of the benchmark dataset supplied, a known limitation for UTKFace's model-estimated age labels in particular, and every reported subgroup below a minimum sample size of \gls{symb:n} $\geq 30$ is flagged as statistically insufficient rather than scored. Full-dataset evaluation is bounded by a documented runtime target rather than guaranteed instantaneous; and if project timeline pressure requires descoping, the platform's own cut-list (\cref{chap:methodology}) defines, in order, which secondary capabilities are dropped first, while the diagnostic core, ingestion, one model interface, the fairness engine, one report format, and the scope-boundary statement itself, is never cut.
=======
Consistent with these boundaries, BiasAperture does not attempt to mitigate bias, retrain or fine-tune models, or generate synthetic demographic data; its function is strictly diagnostic, not corrective. It reports where disparities exist and their statistical strength, but it does not prescribe or apply any remediation. Its findings are only as reliable as the label quality of the benchmark dataset supplied; the current case study uses FairFace, since UTKFace, whose model-estimated age labels were a known limitation, was profiled and cut from the implementation scope, and every reported subgroup below a minimum sample size of \gls{symb:n} \(\geq 30\) is flagged as statistically insufficient rather than scored. Full-dataset evaluation is bounded by a documented runtime target rather than guaranteed instantaneous; and if project timeline pressure requires descoping, the platform's own cut-list (\cref{chap:methodology}) defines, in order, which secondary capabilities are dropped first, while the diagnostic core, ingestion, one model interface, the fairness engine, one report format, and the scope-boundary statement itself, is never cut.
>>>>>>> Stashed changes
